\section{Identification and the periodized basis}
\setcounter{equation}{0}

\subsection*{Proof of Proposition 1}

\emph{Levels and components are determined by the coefficients.} Let $\beta$
and $\tilde\beta$ be measurable maps from $[0,1]^q$ to $\R^p$ with
$\sum_{\cv}\fcoef{\cv}(\bU)X_{\cv} = \sum_{\cv}\tilde\beta_{\cv}(\bU)X_{\cv}$
almost surely, and put $\delta(u) = \beta(u)-\tilde\beta(u)$. Then
$\{\delta(\bU)\T\bX\}^2 = 0$ almost surely, and since $\delta(\bU)$ is
measurable with respect to $\bU$ and $\E\norm{\bX}_2^2 < \infty$,
\[
  0 \;=\; \E\bigl[\{\delta(\bU)\T\bX\}^2 \mid \bU\bigr]
  \;=\; \delta(\bU)\T\,\E(\bX\bX\T\mid\bU)\,\delta(\bU)
  \qquad \text{almost surely.}
\]
The conditional second moment matrix is positive definite almost surely, and a
positive definite quadratic form vanishes only at the origin, so
$\delta(\bU) = 0$ almost surely.

\emph{Levels and components are determined within a coefficient.} Because the
density of $\bU$ is positive Lebesgue almost everywhere, a set is null for the
law of $\bU$ if and only if it is Lebesgue null, so the identity
$\lvl{\cv} + \sum_{\md}\gcomp{\cv}{\md}(u_{\md}) =
\tilde{c}_{\cv} + \sum_{\md}\tilde{g}_{\cv\md}(u_{\md})$ holds Lebesgue almost
everywhere on $[0,1]^q$. Write $\Delta = \lvl{\cv}-\tilde{c}_{\cv}$ and
$\delta_{\md} = \gcomp{\cv}{\md}-\tilde{g}_{\cv\md} \in L_2[0,1] \subset
L_1[0,1]$, so that $\Delta + \sum_{\md}\delta_{\md}(u_{\md}) = 0$ almost
everywhere. Each summand is integrable on the cube by Tonelli's theorem.
Integrating in all coordinates and using $\int_0^1\delta_{\md} = 0$, which
holds because both parametrizations satisfy (2.2), gives $\Delta = 0$; fixing
$\md$ and integrating in the remaining coordinates gives
$\delta_{\md}(u_{\md}) = 0$ for almost every $u_{\md}$. \qed

\begin{remark}
\label{rem:centering}
If the centering $\int_0^1 \gcomp{\cv}{\md} = 0$ is replaced by
$\E\{\gcomp{\cv}{\md}(U_{\md})\} = 0$, the same proof applies under the
two sided density bound of Assumption 3, since then $L_2[0,1]$ and
$L_2(P_{U_{\md}})$ coincide as sets: the step that uses Fubini's theorem
yields $\delta_{\md}$ constant almost everywhere, and the constant is its
mean, which is zero. The two parametrizations are related by
$\lvl{\cv} \mapsto \lvl{\cv} + \sum_{\md}\E\{\gcomp{\cv}{\md}(U_{\md})\}$, as
stated in Section 2.1 of the paper. The difference is not cosmetic for the
estimator: the design of Section 2.3, which omits the constant scaling
function, estimates the version centered with respect to Lebesgue measure.
\end{remark}

\subsection*{Proof of Lemma 1}

\emph{The basis.} The periodization of the scaling function is
$\phi^{\rm per}_{00}(u) = \sum_{\iota}\phi(u+\iota)$, a finite sum at each $u$
because the support is compact, and it equals the constant function one, since
the integer translates of an orthonormal scaling function with
$\int\phi = 1$ form a partition of unity \citep[Section~7.5.1]{Mallat-2009}. For a wavelet, exchanging the finite
sum with the integral,
$\int_0^1\wav{\lev}{\tsl} = \int_{\R}2^{\lev/2}\psi(2^{\lev}y-\tsl)\,dy =
2^{-\lev/2}\int_{\R}\psi = 0$ by the vanishing moment of order zero. That
$\{\mathbf{1}\}\cup\{\wav{\lev}{\tsl}\}$ is an orthonormal basis of $L_2[0,1]$
is the standard construction of the periodized multiresolution analysis
\citep[Section 9.3]{Daubechies-1992}. Since $\dim\VJ = 2^J$ and the $2^J-1$
wavelets of level below $J$ are orthogonal to the constant, $\WJ$ is exactly
the set of elements of $\VJ$ with zero integral.

\emph{Injectivity.} Suppose two parameter vectors give the same regression
function almost surely. The functions
$\beta^J_{\cv}(u) = \lvl{\cv}+\sum_{\md}g^J_{\cv\md}(u_{\md})$, with
$g^J_{\cv\md} = \sum_{\lev\tsl}\coef{\cv}{\md}{\lev}{\tsl}\wav{\lev}{\tsl}$,
are measurable and bounded, so the first part of Proposition 1 gives
$\beta^J_{\cv} = \tilde\beta^J_{\cv}$ almost surely; each $g^J_{\cv\md}$ has
zero integral by the previous paragraph, so the second part gives
$\lvl{\cv} = \tilde{c}_{\cv}$ and
$\sum_{\lev\tsl}(\btheta-\tilde\btheta)_{\cv\md,\lev\tsl}\wav{\lev}{\tsl} = 0$
almost everywhere, whence $\btheta = \tilde\btheta$ by orthonormality. For the
Gram matrix, $\E\norm{\bZ}_2^2 < \infty$ because $\E\norm{\bX}_2^2 < \infty$
and each wavelet is bounded, and for $a \ne 0$,
$a\T\Gram a = \E\{(a\T\bZ)^2\} = \norm{f_a}^2_{L_2(P)} > 0$ by injectivity,
where $f_a$ is the regression function with parameter $a$.

\emph{Nullity when the scaling column is kept.} With $\phi_{00}\equiv 1$ the
column $X_{\cv}\phi_{00}(U_{\md})$ coincides with the unpenalized column
$X_{\cv}$, for each of the $pq$ pairs, so the $pq$ vectors that oppose one to
the other lie in the kernel and are linearly independent. There are no others:
a kernel vector defines
$\sum_{\cv}X_{\cv}\{a_{\cv} + \sum_{\md}(\alpha_{\cv\md} +
\sum_{\lev\tsl}a_{\cv\md,\lev\tsl}\wav{\lev}{\tsl}(U_{\md}))\} = 0$ almost
surely; by Proposition 1 applied to the part with zero integral,
$a_{\cv\md,\lev\tsl} = 0$ and $a_{\cv}+\sum_{\md}\alpha_{\cv\md} = 0$, which
is the space of dimension $pq$ spanned by those vectors. \qed

\begin{remark}[The interval basis]
\label{rem:cdv}
With the basis of \citet*{cohen1993wavelets} the wavelets of level
$\lev \ge \jzero$ are orthogonal to $V_{\jzero}$, which contains the
constants, so they still have zero integral and the penalized part of each
block satisfies (2.2) with no restriction. The $2^{\jzero}$ scaling functions
do not: the constant is the combination $1 = \sum_{\tsl}\mu_{\tsl}
\phi^{\rm int}_{\jzero\tsl}$ with
$\mu_{\tsl} = \int_0^1\phi^{\rm int}_{\jzero\tsl}$, and the subspace with zero
integral is $\{\mu\T\alpha = 0\}$. Keeping the $2^{\jzero}$ scaling columns of
each block unrestricted repeats the constant and makes $\Gram$ singular with
nullity $pq$, as in the last part of Lemma 1; reparametrizing
$\alpha_{\cv\md} = Q\gamma_{\cv\md}$ with $Q$ an orthonormal basis of
$\{\mu\T\alpha = 0\}$ restores injectivity, at the price of
$pq(2^{\jzero}-1)$ unpenalized parameters.
\end{remark}

%%%%%%%%%%%%%%%%%%%%%%%%%%%%%%%%%%%%%%%%%%%%%%%%%%%%%%%%%%%%%%%%%%%%%%%%%%%%%%
% derivations/02-aproximacao-besov.tex
